\section{Motivation and Design Objectives}
\label{sec:heuristics_motivation}
\hspace{\parindent}
The motivation for developing prioritization heuristics arises from the fact that the same planning instance may produce substantially different results depending only on the order in which robots are planned. In constrained environments, assigning priority to one robot may reduce its own planning difficulty while increasing the constraints imposed on the remaining robots.

Since the most suitable ordering depends on the spatial distribution of robots, their goals, the structure of the roadmap, and the interactions between paths, a single fixed prioritization rule is unlikely to perform consistently across all scenarios. For this reason, the objective of this work is not to identify one dominant heuristic, but to design a diverse set of complementary ordering strategies. These heuristics should produce priority orderings that are informative enough to improve planning feasibility or solution quality, while remaining computationally lightweight when compared with the execution time of TEA*. This cost requirement is particularly important in the context of online fleet coordination. Since prioritization is performed before planner execution, excessive heuristic computation time could reduce or eliminate the practical benefit obtained from improved orderings.

\section{Heuristic Design Methodology}
\label{sec:heuristic_design_methodology}
\hspace{\parindent}
The heuristics developed in this work follow a common score-based formulation. Given a set of robots

\begin{equation}
    \mathcal{R} = \{r_1, r_2, \ldots, r_n\}
\end{equation}
, each heuristic assigns a priority score, $score_i$, to every robot, $r_i$, according to the information available in the current planning instance. For a given heuristic, this score is denoted as:

\begin{equation}
    score_i = h(r_i, \mathcal{P})
\end{equation}
, where $\mathcal{P}$ represents the current planning instance, including information such as the robots' current positions, target vertices, roadmap graph properties, estimated paths, predicted conflicts, or information from previous coordination cycles.

The final priority ordering is obtained by sorting the robots according to their scores. Depending on the heuristic, robots with higher scores may be planned first or last.
This formulation provides a common implementation interface for all proposed heuristics. Each heuristic differs in the information used to compute $score_i$ and in the direction of the sorting operation. As a result, different prioritization strategies can be used within the same TEA* planning pipeline and later explored in parallel.

The developed heuristics are organized into five main groups according to the type of information they exploit:

\begin{itemize}
    \item \textbf{Roadmap distance heuristics}, based on distances or estimated path lengths in the roadmap graph;

    \item \textbf{Topological heuristics}, based on structural properties of the graph and its constrained regions;

    \item \textbf{Surroundings-based heuristics}, based on the spatial distribution of robots around their current positions or target goals;

    \item \textbf{Conflict-based heuristics}, based on estimated interactions between tentative robot paths;

    \item \textbf{Replanning heuristics}, based on information from previous plans or coordination cycles.
\end{itemize}

Together, these groups cover complementary dimensions of the coordination problem, allowing the proposed heuristics to address robot prioritization from different perspectives. Each group is examined in the following sections, where the corresponding heuristics are described in terms of their motivation, scoring procedure, and computational implications.

\section{Roadmap Distance Heuristics}
\label{sec:roadmap_distance_heuristics}
\hspace{\parindent}
Roadmap distance heuristics define robot priority orderings using measures related with distance between each robot and its target. The main intuition behind this family is that robots associated with longer movements may have a higher probability of interacting with other robots, since their paths tend to cover a larger portion of the environment and may traverse more shared resources.

These heuristics have the advantage of being computationally lightweight, as the priority permutation is computed from simple distance measures or shortest-path estimates. However, they do not explicitly account for path conflicts, local robot density, or topological constraints. Instead, they rely on travel distance as a simple indicator of the relative difficulty of planning for each robot.

In this work, two criteria based on distance are considered, namely the distance between each robot and its target and the length of an estimated path in the roadmap graph. These criteria give rise to the Furthest Robot and Longest Path heuristics, respectively.


\subsection{Furthest Robot}
\label{subsec:furthest_robot}
\hspace{\parindent}
The Furthest Robot heuristic assigns priorities according to the direct distance between each robot and its target. 

For each robot $r_i$, let $s_i = (x_i, y_i)$ denote its current position and let $g_i = (x_{g_i}, y_{g_i})$ denote its target position. The score assigned to robot $r_i$ is computed as the Euclidean distance between these two points:

\begin{equation}
    score_i = dist(s_i, g_i)
    =
    \sqrt{(x_{g_i} - x_i)^2 + (y_{g_i} - y_i)^2}.
\end{equation}

Figure~\ref{fig:furthest_robot_heuristic} illustrates this computation for a single robot-target pair. In the example, the score assigned to robot $r_1$ corresponds to the direct distance between its current position, $s_1$, and its target position, $g_1$.

\begin{figure}[H]
    \centering
    \includegraphics[width=0.5\textwidth]{furthest.png}
    \caption{Score computation for the Furthest Robot heuristic.}
    \label{fig:furthest_robot_heuristic}
\end{figure}



\subsubsection{Ordering variants}
\label{subsubsec:furthest_robot_ordering_variants}

Two ordering variants were considered. In the \textit{Furthest Robot First} variant, robots with higher distance scores are planned earlier, following the intuition that longer movements should be planned before the search space becomes more constrained.

In the \textit{Furthest Robot Last} variant, the ordering is reversed, and robots with lower distance scores are planned earlier. This strategy may be useful when robots close to their targets can complete their movements quickly, reducing the number of active robots and simplifying the remaining coordination problem.

\subsubsection{Pseudocode}
\label{subsubsec:furthest_robot_pseudocode}

The procedure is summarized in Algorithm~\ref{alg:furthest_robot}.

\begin{algorithm}[H]
\caption{Furthest Robot Heuristic}
\label{alg:furthest_robot}
\begin{algorithmic}[1]
\Require Set of robots $\mathcal{R}$, current positions $s_i$, target positions $g_i$, ordering mode $mode$
\Ensure Priority ordering $\pi$

\For{each robot $r_i \in \mathcal{R}$}
    \State $score_i \gets dist(s_i, g_i)$
\EndFor

\If{$mode = \textit{First}$}
    \State $\pi \gets$ robots sorted by decreasing $score_i$

\ElsIf{$mode = \textit{Last}$}
    \State $\pi \gets$ robots sorted by increasing $score_i$
\EndIf

\State \Return $\pi$
\end{algorithmic}
\end{algorithm}
\subsubsection{Computational complexity}
\label{subsubsec:furthest_robot_computational_complexity}

For each robot, the heuristic computes a single Euclidean distance to its target, which requires constant time. Therefore, computing all scores has complexity $O(n)$, where $n$ is the number of robots. The final sorting step dominates the overall cost, resulting in a total complexity of $O(n \log n)$.


\subsection{Longest Path}
\label{subsec:longest_path}


\newpage
\section{Topological Heuristics}
\label{sec:topological_heuristics}

\section{Surroundings-based Heuristics}
\label{sec:surroundings_heuristics}

\section{Conflicts-based Heuristics}
\label{sec:conflicts_heuristics}

\section{Replanning Heuristics}
\label{sec:replanning_heuristics}
